\section{Activations、FFN 与 parallel blocks}

本节把 ReLU/GeLU/GLU/SwiGLU 放在一个门控 FFN 框架下，并解释 parallel block 更偏系统选择。


\begin{knowledgebox}{术语消化：ReLU、GeLU、GLU、SwiGLU}
ReLU 是硬截断非线性，GeLU 用高斯 CDF 做平滑门控。GLU 系列把 FFN 拆成 value 分支和 gate 分支，用逐元素乘法控制哪些 hidden features 通过。SwiGLU 用 Swish 作为 gate，现代 LLaMA/PaLM 风格模型普遍采用。
\end{knowledgebox}

\begin{importantbox}{FFN 公式：普通激活与门控激活}
普通 FFN 可以写成
\[
\mathrm{FFN}(x)=\phi(xW_1)W_2.
\]
GLU/SwiGLU 类门控 FFN 则写成
\[
\mathrm{SwiGLU}(x)=\bigl(xW_v\bigr)\odot \mathrm{Swish}(xW_g)\,W_o,
\]
其中 \(W_v\) 是 value projection，\(W_g\) 是 gate projection，\(\odot\) 是逐元素乘法。门控结构的代价是多一支投影，所以很多模型会把中间维度缩到约 \(2/3\)，让参数量和 FLOPs 与普通 FFN 可比。
\end{importantbox}

\begin{knowledgebox}{符号说明：FFN 与 RoPE 公式}
\(x\) 表示当前位置 token 的 hidden state；\(W_1,W_2,W_v,W_g,W_o\) 是可学习投影矩阵；\(\phi\) 是 ReLU/GeLU/Swish 等非线性；\(\odot\) 表示逐元素乘法。RoPE 中 \(i\) 是 token 位置，\(k\) 是第 \(k\) 个二维旋转平面，\(\theta_k\) 是该平面的旋转频率。公式说明的是机制形状，不代表唯一实现细节。
\end{knowledgebox}


这一部分先故意展示 activation zoo，因为名称最容易制造“每个函数都完全不同”的错觉。更好的分类方式只有两轴：第一，普通 FFN 还是 gated FFN；第二，非线性使用 ReLU、GeLU、Swish 还是近似线性。这样 GeGLU、ReGLU、SwiGLU 与 LiGLU 都只是同一门控骨架上的 gate choice，而不是四套无关架构。

\begin{figure}[H]
\centering
\includegraphics[width=0.88\textwidth]{slide-019.jpg}
\caption{Slide 20：activation zoo；大量名称需要按普通 FFN 与 gated FFN 重新分类。}
\figsource{CS336 Spring 2026 Lecture 3 官方幻灯片。}
\end{figure}

\begin{knowledgebox}{术语表不要只写“更平滑”}
ReLU 的核心是零阈值截断；GeLU 用输入在高斯分布下的概率做连续门控；Swish 近似为输入乘 sigmoid；GLU 家族则显式增加一条 gate projection。比较时必须同时计算参数量、两支投影的中间维度和可用 kernel，否则“激活函数更好”可能只是模型更大或实现更快。
\end{knowledgebox}

Slide 21 把公式和代表模型放在一起，展示了明显的历史迁移：早期 Transformer、T5、Gopher 与 Chinchilla 常见 ReLU；GPT 系列与 BLOOM 广泛使用 GeLU；2023 年后的 LLaMA、PaLM 与许多新模型则更常用 SwiGLU 或 GeGLU。这个时间趋势提供 strong prior，但最终因果证据还要看后续的等参数消融。

\begin{figure}[H]
\centering
\includegraphics[width=0.88\textwidth]{slide-020.jpg}
\caption{Slide 21：ReLU、GeLU 与 SwiGLU/GeGLU 的公式和代表模型。}
\figsource{CS336 Spring 2026 Lecture 3 官方幻灯片。}
\end{figure}

\begin{warningbox}{模型名单不是公平对照}
不同模型同时改变数据、规模、optimizer 与训练 token，不能从“新模型常用 SwiGLU”直接推出性能增益。名单只说明可用性和社区共识；真正判断 gate 是否值得，需要在相同参数预算、相近 FLOPs 与相同训练 recipe 下比较。
\end{warningbox}


\begin{figure}[H]
\centering
\includegraphics[width=0.88\textwidth]{slide-021.jpg}
\caption{Slide 22: Gated activations}
\figsource{CS336 Spring 2026 Lecture 3 官方幻灯片。}
\end{figure}

\noindent\textbf{展开说明：}本页解释 GLU 的机制：一支产生候选值，另一支产生 gate，用逐元素乘法调制 FFN 输出。


\begin{figure}[H]
\centering
\includegraphics[width=0.88\textwidth]{slide-022.jpg}
\caption{Slide 23: Gated FF variants}
\figsource{CS336 Spring 2026 Lecture 3 官方幻灯片。}
\end{figure}

\noindent\textbf{展开说明：}本页比较 GeGLU、ReGLU、SwiGLU、LiGLU 等变体，核心差异在 gate 分支用什么非线性。


\begin{figure}[H]
\centering
\includegraphics[width=0.88\textwidth]{slide-023.jpg}
\caption{Slide 24: Do GLUs work}
\figsource{CS336 Spring 2026 Lecture 3 官方幻灯片。}
\end{figure}

\noindent\textbf{展开说明：}本页展示 Shazeer 2020 的证据：gated variants 通常稳定优于普通 FFN。

\begin{knowledgebox}{读图：Slide 24 应该怎么看}
这页不是装饰性截图，而是本节论点的证据页。先看图中模型/论文/曲线或表格的比较对象，再看它们支持的趋势：哪些设计已经形成共识，哪些只是局部实验结果，哪些主要由运行时或推理成本推动。读这类页时要区分 quality evidence、runtime evidence 和 stability evidence，不能把一种证据直接替换成另一种。
\end{knowledgebox}



Shazeer 2020 的单篇结果还不足以形成默认值，因此 Slide 25 加入 Narang et al. 的独立结果。读这页时先确认比较是否保持参数量或 compute 接近，再看不同任务和规模上 gated variants 的方向是否一致。多项工作重复观察到收益，才把“有趣变体”提升为现代 recipe 的高概率起点。

\begin{figure}[H]
\centering
\includegraphics[width=0.88\textwidth]{slide-024.jpg}
\caption{Slide 25：独立工作对 gated activation 收益的复核。}
\figsource{CS336 Spring 2026 Lecture 3 官方幻灯片；Narang et al. 2020。}
\end{figure}

\begin{knowledgebox}{读图：重复证据比单个最好点更重要}
不要只找某一列的最小 loss，而应看不同设置下 gated variants 是否多数优于普通 ReLU/GeLU，以及优势是否需要额外参数才能出现。跨论文方向一致说明机制具有一定稳健性；具体选择 SwiGLU 还是 GeGLU，仍可能由 kernel、数值稳定和现有实现决定。
\end{knowledgebox}

阶段性 recap 进一步给出保守结论：普通 GeLU/ReLU 仍能训练出强模型，GLU 不是“模型能否工作”的必要条件；但在等预算证据和大量现代模型实践下，SwiGLU/GeGLU 已成为更好的默认实验起点。Nemotron 的 squared ReLU 等例外值得记录，却不能推翻主流证据，只说明搜索空间仍未关闭。

\begin{figure}[H]
\centering
\includegraphics[width=0.88\textwidth]{slide-025.jpg}
\caption{Slide 26：activation 与 gating 的保守总结及少数例外。}
\figsource{CS336 Spring 2026 Lecture 3 官方幻灯片。}
\end{figure}

\begin{importantbox}{选择 activation 的实际流程}
先用实现成熟的 SwiGLU 或 GeGLU 作为 baseline，并把中间维度按参数预算缩放；再用小规模 ablation 比较 validation loss、step time、peak memory 和数值稳定性。只有当特殊 activation 在相同预算下持续改进这些指标，才值得承担非标准 kernel 与部署复杂度。
\end{importantbox}


\begin{figure}[H]
\centering
\includegraphics[width=0.88\textwidth]{slide-026.jpg}
\caption{Slide 27: Serial vs parallel layers}
\figsource{CS336 Spring 2026 Lecture 3 官方幻灯片。}
\end{figure}

\noindent\textbf{展开说明：}本页比较标准串行 block：先 attention 后 MLP。它引出 parallel block 作为系统/延迟选择。



标准 block 串行执行 attention 与 MLP，而 Slide 28 展示 GPT-J、PaLM、GPT-NeoX 等采用 parallel layers：两个 branch 读取同一个 normalized residual state，再把输出一起加回。它减少依赖链，并可能共享 LayerNorm、并行调度或融合矩阵乘；代价是计算图和串行 block 不完全等价，质量需要单独验证。

\begin{figure}[H]
\centering
\includegraphics[width=0.88\textwidth]{slide-027.jpg}
\caption{Slide 28：parallel layers 的计算图、代表模型与系统动机。}
\figsource{CS336 Spring 2026 Lecture 3 官方幻灯片。}
\end{figure}

\begin{knowledgebox}{Parallel block 主要解决什么}
这项选择的第一动机通常是 latency 与 kernel scheduling，而不是提出更强的非线性。若两个 branch 可以同时执行，或把投影拼成更大的矩阵乘，就可能提高硬件利用率；但端到端收益取决于并行 runtime、通信和 batch shape。没有实际 profiler 数据时，不能只凭结构图断言更快。
\end{knowledgebox}


\begin{figure}[H]
\centering
\includegraphics[width=0.88\textwidth]{slide-028.jpg}
\caption{Slide 29: Architecture summary}
\figsource{CS336 Spring 2026 Lecture 3 官方幻灯片。}
\end{figure}

\noindent\textbf{展开说明：}本页总结 architecture variations：pre-norm、RMSNorm、SwiGLU、parallel block 等哪些是强共识，哪些仍是局部选择。

\subsection{本章小结}
本节的关键不是记住所有模型名，而是把公开模型的设计选择转化成可迁移判断：共识默认值可以作为起点，例外需要知道原因，涉及系统和推理成本的选择必须结合资源账本讨论。

